\documentclass[12pt]{article}

\usepackage{amsmath}
\usepackage{geometry}
\usepackage{hyperref}
\usepackage{parskip}

\geometry{a4paper, margin=1in}

\title{A Latent Log-Ratio Model for Practical Shooting Ratings}
\author{}
\date{}

\begin{document}
	
	\maketitle
	
	\begin{abstract}
		Traditional rating systems like Elo and Glicko were designed for zero-sum, pairwise competitions (e.g., chess). Applying them to practical shooting sports (USPSA, IDPA, ICORE) introduces severe structural friction. Practical shooting is massively multiplayer, with matches featuring dozens to hundreds of competitors simultaneously. Furthermore, it is continuous and non-zero-sum: one competitor's high performance does not come at the cost of another's, and scores are typically given as a ratio to the winner.
		
		There are adaptations to Elo and Glicko for continuous, non-zero-sum games, but they are just that: adaptations. All make compromises to shoehorn the continuous nature of the game into the binary, head-to-head model of the rating system.
		
		This paper proposes a Latent Log-Ratio Model, a Bayesian state-space approach that operates directly on ratio-space finishes. The model uses the logarithmic structure of multiplicative performance ratios to eliminate the winner's-curse bias inherent in percentage scoring while preserving globally log-linear rating semantics for prediction. It combines a linear-time match baseline with selective pairwise residual corrections, a robust and weakly regularized leave-one-out baseline estimator, score-dependent tail noise for very weak finishes, a match-context weak-field uncertainty term for tiny bottom-heavy fields, per-competitor behavioral volatility tracked via exponential moving averages, and innovation-based adaptive variance to prevent the monotonic certainty collapse common in basic Kalman filters.
		
		Though designed with practical shooting in mind, the mathematics are not specific to that sport, and the system should function well when applied to any game where a ratio-to-winner finish order can be established.
	\end{abstract}
	
	\section{Key Insights and Concepts}
	
	\subsection{The Multiplicative Nature of Scoring}
	In practical shooting, true skill is fundamentally multiplicative, whether measured by hit factor  in USPSA, or total time in time-plus sports like IDPA and ICORE. If Competitor A takes 10 seconds to shoot a stage and Competitor B takes 20 seconds, Competitor A is twice as fast. On a longer stage, the times might be 20 seconds and 40 seconds; the absolute margin changes, but the $2\times$ ratio remains constant.
	
	\subsection{The ``Winner's Curse'' in Percentage Scoring}
	Matches are scored as percentages of the top competitor's performance. Using these raw percentages injects the winner's variance into the entire field.
	\begin{itemize}
		\item \textbf{Hit Factor} (USPSA, where $P_i$ is a competitor's finish percentage in $(0.0 - 1.0)$): $S_i = \frac{P_i}{1.0}$
		\item \textbf{Time-Plus} (IDPA/ICORE): $S_i = \frac{T_{min}}{T_i}$
	\end{itemize}
	If the match winner performs exceptionally well, every other competitor's percentage drops, mathematically suggesting the field performed poorly when they did not.
	
	\subsection{The Logarithmic Bridge}
	Let $P_i$ denote a competitor's absolute performance on any positive ratio scale where larger is better, and let $P_{max}$ denote the winning performance in the match. The match percentage is then
	\begin{equation}
		S_i = \frac{P_i}{P_{max}}
	\end{equation}
	For any two competitors $A$ and $B$, the difference of their log-percentages is
	\begin{equation}
		\ln(S_A) - \ln(S_B) = \ln\left(\frac{P_A}{P_{max}}\right) - \ln\left(\frac{P_B}{P_{max}}\right)
	\end{equation}
	Using $\ln(x/y) = \ln(x) - \ln(y)$, this becomes
	\begin{equation}
		\ln(S_A) - \ln(S_B) = (\ln(P_A) - \ln(P_{max})) - (\ln(P_B) - \ln(P_{max}))
	\end{equation}
	and the common anchor cancels exactly:
	\begin{equation}
		\ln(S_A) - \ln(S_B) = \ln(P_A) - \ln(P_B) = \ln\left(\frac{P_A}{P_B}\right)
	\end{equation}
	This is the central trick of the model: once the winner anchor cancels, the observed log-ratio of absolute performances can be treated as a direct estimate of latent rating difference,
	\begin{equation}
		\ln\left(\frac{P_A}{P_B}\right) \approx R_A - R_B
	\end{equation}
	Thus, the winner's score appears only as a nuisance normalization constant in the published percentages; it does not survive pairwise comparison in log space. The resulting quantity is the pure ratio of the competitors' absolute performances, independent of who won the match or by how much. In time-plus sports, this same derivation applies by choosing $P_i$ to be any proportional speed variable such as $T_{min}/T_i$, which yields $\ln(S_A) - \ln(S_B) = \ln(T_B/T_A)$.
	
	\subsection{Ratings as Expected Log-Ratios}
	This cancellation property leads directly to the model's core identity. Since ratings $R_i$ represent expected log-performance, the difference between any two ratings is the expected log-ratio of their performances:
	\begin{equation}
		R_i - R_j = E[\ln(S_i) - \ln(S_j)] = E\!\left[\ln\!\left(\frac{S_i}{S_j}\right)\right]
	\end{equation}
	
	Equivalently, the expected ratio of two competitors' performances is:
	\begin{equation}
		E\!\left[\frac{S_i}{S_j}\right] \approx \exp(R_i - R_j)
	\end{equation}
	
	The rating difference \textit{is} the expected performance difference: a rating gap of $\ln(2) \approx 0.69$ means one competitor is expected to perform at twice the level of the other. This relationship remains exact throughout the implemented model. Deep-tail pathologies are handled by inflating observation noise for very weak finishes, not by deforming the log-ratio mean link, so the prediction mapping remains globally reversible through the exponential function.
	
	\section{System Architecture and Algorithm}
	The system models each competitor's latent (hidden) skill as a Gaussian distribution in log-performance space, augmented with a per-competitor volatility estimate.
	
	\subsection{State Representation}
	Each competitor $i$ is defined by three variables:
	\begin{itemize}
		\item \textbf{$R_i$ (Rating):} The expected mean of their natural log performance. An average baseline competitor can be anchored at 0.
		\item \textbf{$V_i$ (Variance):} The uncertainty in the rating. High for new competitors, decreasing as they record more scores, but capable of increasing in response to surprising results.
		\item \textbf{$\sigma_i^2$ (Competitor Volatility):} A running estimate of how erratic the competitor's performances are (behavioral noise for update damping). Tracks whether a competitor is consistent or prone to large swings. It is analogous in role to Glicko-2's volatility, but computed via an exponential moving average rather than iterative root-finding. Win-probability mixture weights still omit $\sigma_i^2$ (not uncertainty about who wins); optional scaled $\kappa\sigma_i^2$ may enter only posterior-predictive bands (see \S\,\ref{sec:prediction}).
	\end{itemize}
	
	\subsection{System Parameters}
	\begin{itemize}
		\item \textbf{$s$, $b$ (Display Scale):} User-facing ratings use an affine map $R^{\mathrm{disp}} = s\cdot R + b$ (defaults are sport-specific in software). Variance and volatility in display space scale as $s^2$ times latent variance; display standard deviations scale as $|s|$ times latent standard deviation.
		\item \textbf{$\sigma_{sport}^2$ (Sport Volatility):} A tuned constant representing the inherent noise of a single match observation in the sport. This is the irreducible variance that even a perfectly consistent competitor experiences due to environmental factors, range conditions, and the stochastic nature of shooting. Tuned once per sport; does not vary per competitor. (This value will ordinarily grow substantially when considering individual stages vs. matches as the rating event.)
		\item \textbf{$\sigma_{drift}^2$ (Skill Drift Rate):} The rate at which rating uncertainty grows during periods of inactivity. Represents the possibility that a competitor's true skill changes when they are not being observed. Applied as $V_i \leftarrow V_i + \sigma_{drift}^2 \times \Delta t$ between matches, where $\Delta t$ is measured in rating periods (then capped at $V_{\max}$ in software; see $V_{\max}$ below).
		\item \textbf{$V_{\mathrm{start}}$ (Starting Variance):} The committed variance assigned to competitors with no rating history (prior width for new shooters).
		\item \textbf{$V_{\max}$ (Maximum Variance):} Upper bound on committed rating variance after Kalman updates and after drift aging. Defaults to $V_{\mathrm{start}}$ for backward compatibility but may be set larger so veterans can carry more epistemic uncertainty without widening the new-shooter prior.
		\item \textbf{$\beta$ (Volatility Adaptation Rate):} Controls how quickly the per-competitor volatility $\sigma_i^2$ adapts to recent results. A value of 0.1--0.15 gives an effective lookback window of approximately 7--10 matches.
		\item \textbf{$\gamma$ (Surprise Adaptation Rate):} Controls how aggressively the variance $V_i$ responds to unexpectedly surprising results. Recommended range: 0.05--0.1.
		\item \textbf{$\alpha$ (Pairwise Blend Weight):} Controls how much structural information from pairwise opponent comparisons is blended into the baseline-derived observed performance. When $\alpha = 0$, the model reduces to the pure baseline model.
		\item \textbf{$s_0$ (Tail Noise Start Percentage):} The finish percentage below which very weak results begin receiving additional observation noise. Above $s_0$, the ordinary variance model is used.
		\item \textbf{$\xi^2$ (Tail Noise Variance):} The maximum additional observation variance applied to the deepest tail. Larger values make the model more skeptical of very weak finishes without changing the log-ratio mean link.
		\item \textbf{$\omega^2$ (Weak-Field Variance):} The maximum additional match-context observation variance applied when a field is both small and bottom-heavy.
		\item \textbf{$n_{max}$ (Weak-Field Max Size):} The field size at or above which weak-field damping decays to zero.
		\item \textbf{$s_w$ (Weak Finish Threshold):} The finish threshold below which a non-winning score counts as ``weak'' for weak-field detection.
		\item \textbf{$q_0$ (Weak Fraction Threshold):} The minimum fraction of weak non-winners required before the weak-field term activates.
		\item \textbf{$c$ (Baseline Robustness Threshold):} A Huber-style threshold, measured in residual standard deviations, used to downweight extreme field anchors when estimating match baseline.
		\item \textbf{$\tau^2$ (Match Difficulty Prior Variance):} A weak prior variance on the match baseline. Its precision $1/\tau^2$ acts as a virtual anchor at zero match difficulty and is best treated as a light regularization term rather than a literal estimate of real-world match-to-match difficulty.
		\item \textbf{$\sigma_{pred}^2$ (Prediction Sport Variance):} Idiosyncratic sport noise used only in the prediction pipeline (win-probability denominators and predictive variance). It is typically set smaller than $\sigma_{sport}^2$ because common-mode match difficulty largely cancels in relative finish ratios; the rating update continues to use the full $\sigma_{sport}^2$.
		\item \textbf{$\kappa$ (Prediction Behavioral Scale):} Dimensionless coefficient in $[0,1]$ (typically $\ll 1$) that scales how much per-competitor behavioral variance $\sigma_i^2$ enters \emph{posterior predictive} variance for finish-ratio bands. The rating update still uses full $\sigma_i^2$ in the effective observation noise; $\kappa$ addresses overlap between epistemic rating uncertainty $V_i$ and EMA-estimated $\sigma_i^2$ when forming bands, not a second independent draw in the filter.
	\end{itemize}
	
	\subsection{Match Processing Loop}
	Matches are processed completely online (sequentially). The core baseline computation is $\mathcal{O}(N)$; the optional pairwise residual correction adds $\mathcal{O}(N \cdot k)$ where $k$ is the number of selected opponents per competitor (typically 20--30), yielding total complexity of $\mathcal{O}(N \cdot k)$.
	
	\textbf{Step 1: Calculate Log-Scores}\\
	Floor match percentages at a minimal value $S_{floor}$ (e.g., 0.01) to prevent undefined logarithms, then convert to log space:
	\begin{equation}
		L_i = \ln(\max(S_{floor}, S_i))
	\end{equation}
	To preserve global log-linearity while reducing overconfidence in extremely weak finishes, the implementation keeps this mean link exactly intact and instead introduces score-dependent tail noise in the observation variance.
	
	\textbf{Step 2: Calculate Competitor Weights}\\
	Determine how much influence a competitor should have on establishing the match's difficulty. High-variance and high-volatility competitors have less influence. Very weak finishes can also receive additional observation variance:
	\begin{equation}
		\eta(S_i) =
		\begin{cases}
			0, & S_i \ge s_0 \\
			\xi^2 \left(\frac{s_0 - S_i}{s_0}\right)^2, & S_i < s_0
		\end{cases}
	\end{equation}
	\begin{equation}
		W_i = \frac{1}{V_i + \sigma_{sport}^2 + \sigma_i^2 + \eta(S_i)}
	\end{equation}
	This construction is deliberately asymmetric between mean and variance. The observed score remains
	\begin{equation}
		L_i = \ln(\max(S_{floor}, S_i))
	\end{equation}
	so the globally log-linear interpretation of rating differences is preserved exactly. Tail handling enters only through uncertainty: as a finish moves deeper into the weak tail, the model treats it as less precise evidence about latent skill rather than as evidence with a different mean.
	
	\subsubsection*{Why Tail Noise Is Needed}
	In many practical-shooting datasets, very weak finishes are not reliable linear evidence about top-end skill separation. A competitor finishing at 20\% of the winner is certainly worse than one finishing at 40\%, but it does not follow that the 20\% finisher provides proportionally trustworthy information about how far the winner sits above the rest of the field in latent skill space. Deep-tail results are often contaminated by penalties, stage failures, equipment problems, or simply the nonlinear geometry of ratio scoring in the far tail.
	
	A naive fix would be to compress the log-score itself with a nonlinear map. That can work as a robust update heuristic, but it destroys the model's clean predictive semantics: once the mean link is no longer $\ln(S_i)$, rating differences are no longer trivially interpretable as log-performance ratios. The tail-noise approach instead keeps the central identity intact and expresses skepticism in the probabilistically appropriate place, namely the observation variance.
	
	The piecewise definition of $\eta(S_i)$ encodes two desiderata:
	\begin{enumerate}
		\item For ordinary finishes ($S_i \ge s_0$), the model should reduce to the standard latent log-ratio update, so $\eta(S_i) = 0$.
		\item For weak finishes ($S_i < s_0$), the extra variance should increase smoothly as the score deteriorates, so that extremely weak results have less leverage in baseline anchoring, pairwise residual averaging, and the final Kalman update.
	\end{enumerate}
	
	The quadratic form
	\begin{equation}
		\eta(S_i) = \xi^2 \left(\frac{s_0 - S_i}{s_0}\right)^2
	\end{equation}
	is not the only possible choice, but it has several practical advantages: it is continuous at the threshold, it is monotone in tail depth, it is cheap to compute, and it avoids introducing a second kink in the update geometry. Most importantly, it leaves the mean observation untouched while making low-end results progressively less influential as evidence.
	
	\textbf{Step 3: Establish the Match Baseline ($B$)}\\
	Calculate per-competitor baseline residuals by comparing the field's expected ratings against their log-scores:
	\begin{equation}
		a_i = R_i - L_i
	\end{equation}
	A raw baseline estimate is then:
	\begin{equation}
		B_{raw} = \frac{\sum_{i=1}^N W_i a_i}{\sum_{i=1}^N W_i + \tau^{-2}}
	\end{equation}
	The implementation then robustifies the field anchor with a Huber-style taper. Standardize each residual relative to the raw center:
	\begin{equation}
		z_i = \frac{a_i - B_{raw}}{\sqrt{1/W_i}}
	\end{equation}
	and define robustified baseline weights
	\begin{equation}
		\tilde{W}_i = W_i \cdot \min\!\left(1, \frac{c}{|z_i|}\right)
	\end{equation}
	with the convention that the multiplier is 1 when $z_i = 0$ or $c \le 0$. The robust baseline summary is then
	\begin{equation}
		B = \frac{\sum_{i=1}^N \tilde{W}_i a_i}{\sum_{i=1}^N \tilde{W}_i}
	\end{equation}
	\textit{Note: $B$ acts as the mathematical anchor, answering: ``How difficult was this match compared to the known skill level of the field?''}
	
	\textbf{Step 4: Calculate Pairwise Residuals ($D_i$)}\\
	The match baseline captures the average difficulty of the field, but discards structural information about \textit{whom} each competitor outperformed. In a match with widely varying skill levels, this information is valuable: beating someone rated far above you is a stronger signal than beating someone rated far below you.
	
	For each competitor $i$, the implementation unions several heuristics: the top 10\% of the field by match finish (30\% when the field has fewer than 10 competitors), the top 10\% by rating, and competitors whose finish ratios lie within a multiplicative band around $i$'s ratio (with a wider band in very small fields). For each selected opponent $j$, compute the pairwise residual in log space:
	\begin{equation}
		\delta_{ij} = (L_i - L_j) - (R_i - R_j)
	\end{equation}
	This measures how much better or worse competitor $i$ performed against opponent $j$ relative to what their ratings predicted. The winner's curse is eliminated by the underlying log-ratio geometry, while the score-dependent tail-noise term reduces the leverage of extremely weak finishes through the weights rather than by altering the mean link.
	
	Compute a precision-weighted average of the pairwise residuals:
	\begin{equation}
		D_i = \frac{\sum_{j \in \mathcal{S}_i} w_j \cdot \delta_{ij}}{\sum_{j \in \mathcal{S}_i} w_j}, \quad w_j = \frac{1}{V_j + \sigma_{sport}^2 + \sigma_j^2 + \eta(S_j)}
	\end{equation}
	where $\mathcal{S}_i$ is the selected opponent set for competitor $i$.
	
	\textbf{Step 5: Calculate Observed Performance ($P_i$)}\\
	Blend the baseline-derived performance with the pairwise residual correction. Use the leave-one-out baseline $B_{-i}$ (see \S\,``Propagating Estimation Uncertainty'') to avoid circularity:
	\begin{equation}
		P_i = L_i + B_{-i} + \alpha_i \cdot D_i
	\end{equation}
	When $\alpha = 0$, this reduces to the pure baseline model. As $\alpha$ increases, the structural information about specific opponent matchups contributes more to the observed performance estimate. In software, the nominal blend weight $\alpha$ is multiplied by a taper in $\alpha_i$ that ramps from 0 to 1 as the number of pairwise opponents increases toward 10, so sparse opponent sets rely more on the baseline.
	
	\subsection{The Update Rule}
	For each competitor, update their rating, variance, and volatility.
	
	\textbf{Step 1: Calculate the Innovation}\\
	The innovation measures how surprising the competitor's observed performance was relative to their current rating:
	\begin{equation}
		e_i = P_i - R_i
	\end{equation}
	
	\textbf{Step 2: Calculate the Kalman Gain ($K_i$)}\\
	The Kalman gain balances how much to trust the new observation versus the prior rating. The competitor's own volatility $\sigma_i^2$ enters as additional observation noise, so that erratic competitors receive smaller per-match updates. The schematic form is
	\begin{equation}
		K_i = \frac{V_i}{V_i + \sigma_{sport}^2 + \sigma_i^2},
	\end{equation}
	but in software $\sigma_{sport}^2 + \sigma_i^2$ is replaced by the full effective observation noise $R_{\mathrm{obs},i}$ from \S\,``Propagating Estimation Uncertainty''.
	
	\textbf{Step 3: Update Rating}
	\begin{equation}
		R_i \leftarrow R_i + K_i \cdot e_i
	\end{equation}
	
	\textbf{Step 4: Update Variance (Innovation-Based Adaptive Kalman)}\\
	A standard Kalman filter variance update $V_i \leftarrow V_i(1 - K_i)$ monotonically decreases $V_i$. Over time, this causes the system to become excessively certain about every competitor's rating, regardless of whether recent results support that certainty. To prevent this, we augment the standard update with an innovation-based correction. Let $R_{\mathrm{obs},i}$ denote the \emph{full} effective observation noise from \S\,``Propagating Estimation Uncertainty'' (including baseline and pairwise estimation uncertainty, tail noise, and weak-field terms). After the Student-$t$ robust weighting is applied to $e_i$, the implementation uses
	\begin{equation}
		V_i \leftarrow \min\!\left(V_{\max}, \; V_i(1 - K_i) + \gamma \cdot \max\!\left(0, \; e_i^2 - (V_i + R_{\mathrm{obs},i})\right)\right)
	\end{equation}
	where $V_{\max}$ is the configured maximum committed variance (not necessarily equal to $V_{\mathrm{start}}$).
	
	The $\max$ term activates only when the squared innovation exceeds the expected innovation variance. In plain terms: if a competitor's result is more surprising than the model predicted it should be, the variance increases to reflect reduced certainty. If the result is within expectations, the variance decreases normally.
	
	This is a well-known extension of the Kalman filter (innovation-based adaptive filtering). It provides the same qualitative behavior as Glicko-2's volatility-driven RD adaptation---surprising results make the system more responsive, consistent results make it more stable---but requires no iterative root-finding.
	
	Note that the interaction with $\sigma_i^2$ is self-consistent: erratic competitors (high $\sigma_i^2$) have a higher expected innovation variance, so it takes a genuinely extreme result to trigger a $V_i$ increase. Consistent competitors (low $\sigma_i^2$) are more sensitive to moderate surprises.
	
	\textbf{Step 5: Update Competitor Volatility}\\
	Track how erratic each competitor's performances are using an exponential moving average of squared innovations:
	\begin{equation}
		\sigma_i^2 \leftarrow (1 - \beta) \cdot \sigma_i^2 + \beta \cdot e_i^2
	\end{equation}
	
	This adapts in $\mathcal{O}(1)$ per update: one multiply-add operation. A competitor who consistently performs near their rating will have low $\sigma_i^2$; a competitor who oscillates between brilliant and terrible performances will have high $\sigma_i^2$. The system naturally adjusts: consistent competitors get more responsive updates (higher $K_i$), erratic competitors get more conservative updates (lower $K_i$).
	
	Clamp $\sigma_i^2$ within $[0.01\,V_{\max},\, V_{\max}]$. The ceiling tracks the same scale as the cap on committed rating variance $V_i$; it is no longer tied to $V_{\mathrm{start}}$, so a tight prior for new competitors does not automatically shrink the allowable behavioral-volatility range for experienced shooters. This decouples behavioral volatility limits from $\sigma_{sport}^2$: tightening sport noise for the Kalman gain no longer shrinks the allowable range for $\sigma_i^2$. The floor keeps a small minimum behavioral noise term.
	
	New competitors initialize $\sigma_{i,\mathrm{init}}^2$ to the geometric mean of those endpoints,
	\begin{equation}
		\sigma_{i,\mathrm{init}}^2 = \sqrt{(0.05\,V_{\mathrm{start}})\,V_{\mathrm{start}}} = \sqrt{0.05}\,V_{\mathrm{start}}.
	\end{equation}
	Thus behavioral noise starts mid-band on a multiplicative scale rather than at $\sigma_{sport}^2$, which would double-count uncertainty already carried by $V_{\mathrm{start}}$ and observation noise for shooters with no history.
	
	\textbf{Step 6: Apply Aging Between Matches}\\
	Between rating periods, increase uncertainty to reflect the possibility that true skill has drifted:
	\begin{equation}
		V_i \leftarrow V_i + \sigma_{drift}^2 \times \Delta t
	\end{equation}
	where $\Delta t$ is the elapsed time in rating periods since the competitor's last match. The implementation stores committed $V_i$ capped at $V_{\max}$ and applies the same cap when computing time-aged variance for display and predictions.
	
	\subsection{Propagating Estimation Uncertainty}
	The baseline $B$ and pairwise residual $D_i$ from the Match Processing Loop are \emph{estimates}, not known quantities. In the basic algorithm above, the Kalman gain treats the observation noise as $\sigma_{sport}^2 + \sigma_i^2$, implicitly assuming that $B$ and $D_i$ are measured perfectly. In reality, both carry estimation variance that depends on the field: a baseline anchored by 50 well-known competitors is far more precise than one anchored by 2 newcomers. Ignoring this causes systematic over-updating in small or uncertain fields.
	
	\subsubsection*{Leave-One-Out Baseline}
	To avoid circularity (competitor $i$ influencing their own baseline), replace $B$ with a leave-one-out estimate formed from the robustified weights:
	\begin{equation}
		B_{-i} = \frac{\sum_{j \neq i} \tilde{W}_j a_j}{\sum_{j \neq i} \tilde{W}_j + \tau^{-2}}
	\end{equation}
	
	The variance of this estimate under the precision-weighted mean with prior regularization is:
	\begin{equation}
		\mathrm{Var}(B_{-i}) = \frac{1}{\sum_{j \neq i} \tilde{W}_j + \tau^{-2}}
	\end{equation}
	
	In large fields with well-known competitors, $\sum_{j \neq i} \tilde{W}_j$ is large and the prior is negligible, so $\mathrm{Var}(B_{-i})$ remains tiny. In small or structurally weird fields, robust downweighting and the weak prior both act as stabilizers: extreme anchors lose leverage, and the baseline is gently shrunk toward average difficulty rather than being allowed to run away.
	
	\subsubsection*{Pairwise Estimation Variance}
	The pairwise residual $D_i$ is itself a precision-weighted mean over the opponent set $\mathcal{S}_i$. Its estimation variance is:
	\begin{equation}
		\mathrm{Var}(D_i) = \frac{1}{\sum_{j \in \mathcal{S}_i} w_j}
	\end{equation}
	
	Since $D_i$ enters the observed performance scaled by the effective blend weight $\alpha_i$ (including the taper toward small opponent sets), its contribution to the total observation variance is $\alpha_i^2 \cdot \mathrm{Var}(D_i)$. Where earlier formulas write $\alpha$, read $\alpha_i$ for the tapered value used in code.

	\subsubsection*{Weak-Field Match-Context Damping}
	Some pathological cases are not well-described by per-competitor tail noise alone: tiny matches in which a large share of the non-winners finish extremely far behind the winner. In such fields, the issue is not merely that one weak finish is noisy, but that the match as a whole provides little trustworthy information about top-end skill separation. To model this, introduce a match-level variance term that depends on field size and weak-field composition.
	
	Let $N$ be the field size and define a small-field taper
	\begin{equation}
		f_{size}(N) = \max\!\left(0,\frac{n_{max} - N}{n_{max} - 2}\right)
	\end{equation}
	so that two-person fields receive the full size penalty and fields with $N \ge n_{max}$ receive none.
	
	Let $\mathcal{W}$ be the set of non-winners whose finish percentages fall below the weak threshold:
	\begin{equation}
		\mathcal{W} = \{j : \text{$j$ is not the winner and } S_j < s_w\}
	\end{equation}
	and define the weak non-winner fraction
	\begin{equation}
		q = \frac{|\mathcal{W}|}{N - 1}
	\end{equation}
	If $q < q_0$, no weak-field penalty is applied. Otherwise, define a severity term by measuring how far the weak finishes extend below the threshold in log space:
	\begin{equation}
		g = \min\!\left(1,\frac{1}{|\mathcal{W}| \cdot (-\ln s_w)} \sum_{j \in \mathcal{W}} \max\!\left(0,\ln\!\left(\frac{s_w}{\max(S_{floor}, S_j)}\right)\right)\right)
	\end{equation}
	The match-level pathology score is then
	\begin{equation}
		p =
		\begin{cases}
			0, & q < q_0 \\
			f_{size}(N) \cdot q \cdot g, & q \ge q_0
		\end{cases}
	\end{equation}
	and the resulting weak-field observation variance is
	\begin{equation}
		\sigma_{path}^2 = \omega^2 \cdot p
	\end{equation}
	This term is shared across competitors in the same match. Conceptually, it says: ``this field may be real, but it is not an especially precise instrument for measuring elite-vs-elite skill.'' Current USPSA-oriented implementation defaults are $s_0 = 0.60$, $\xi^2 = 0.06$, $\omega^2 = 0.50$, $n_{max} = 10$, $s_w = 0.60$, and $q_0 = 0.40$.
	
	\subsubsection*{Effective Observation Noise}
	Combining all variance sources, the total observation noise for competitor $i$'s performance estimate is:
	\begin{equation}
		R_{obs,i} = \sigma_{sport}^2 + \sigma_i^2 + \eta(S_i) + \mathrm{Var}(B_{-i}) + \alpha_i^2 \cdot \mathrm{Var}(D_i) + \sigma_{path}^2
	\end{equation}
	
	This replaces $\sigma_{sport}^2 + \sigma_i^2$ in the Kalman gain (Step~2 of the Update Rule), the innovation-based variance update (Step~4), and the $z$-score used for robust weighting:
	\begin{equation}
		K_i = \frac{V_i}{V_i + R_{obs,i}}, \qquad z_i = \frac{e_i}{\sqrt{V_i + R_{obs,i}}}
	\end{equation}
	
	The effect is automatic and principled: large fields with strong opponents and non-pathological finishes produce small $R_{obs,i}$ (close to the original $\sigma_{sport}^2 + \sigma_i^2$), so rating updates proceed normally. Small fields, uncertain opponents, extremely weak individual finishes, and tiny bottom-heavy fields inflate $R_{obs,i}$, which simultaneously reduces $K_i$ (smaller updates) and reduces $|z_i|$ (making the Student-$t$ robust weight less aggressive). No deformation of the log-ratio mean link is required.
	
	\subsubsection*{Robust Outlier Weighting}
	Even with correct observation noise, individual match results can be corrupted by equipment failure, rule violations, or other non-skill factors. A Student-$t$ weight on the innovation provides a final layer of robustness:
	\begin{equation}
		w_t(z_i) = \frac{\nu + 1}{\nu + z_i^2}
	\end{equation}
	where $\nu$ scales with field size (e.g., $\nu = \mathrm{clamp}(2 + 0.5N, \; 2, \; 30)$). The weighted innovation $e_i \leftarrow w_t(z_i) \cdot e_i$ is used in all subsequent update steps. Because $z_i$ is now computed against the full $R_{obs,i}$, the threshold for ``surprising'' is correctly calibrated: a $z$-score of 3 in a 100-person match genuinely indicates an outlier, whereas the same raw innovation in a 2-person match would produce a much smaller $z$-score (due to the large $\mathrm{Var}(B_{-i})$ term), and the Student-$t$ weight would remain close to 1.
	
	\subsection{Handling Non-Ratio Anomalies: Time Bonuses and Micro-Stages}
	
	The Latent Log-Ratio model fundamentally assumes that the scoring data exists on a \textbf{Ratio Scale} (where zero represents an absolute absence of time, and ratios between values are physically meaningful). However, specific rule sets can violate this assumption:
	
	\begin{enumerate}
		\item \textbf{Time Bonuses (ICORE):} Rules that subtract fixed time for accuracy can result in near-zero or even negative stage times. This converts the data into an Interval Scale, where log-ratios are mathematically undefined or heavily distorted.
		\item \textbf{Micro-Stages (Speed Shoots):} Even without bonuses, very fast stages (e.g., a 2.0-second stage) cause artificial volatility magnification. A 0.5-second stumble on a 2.0-second stage yields a massive 25\% ratio penalty, whereas a 0.5-second stumble on a 20-second stage is mathematically negligible. 
	\end{enumerate}
	
	\subsubsection*{Additive Time Normalization (Time Padding)}
	To restore the ratio-scale properties and prevent infinite variance from undefined logarithms, we apply an additive normalization constant ($k$) to all raw times on stages that fall below a minimum time threshold ($T_{floor}$). 
	
	Let $T_{min}$ be the winning (lowest) time on a given stage. If $T_{min} < T_{floor}$ (e.g., $T_{floor} = 10.0$ seconds), we calculate a padding constant:
	\begin{equation}
		k = T_{floor} - T_{min}
	\end{equation}
	
	This constant is then added to every competitor's time on that stage before converting the finishes into percentages:
	\begin{equation}
		T_i' = T_i + k
	\end{equation}
	
	This translation guarantees that all $T_i' \geq T_{floor} > 0$, ensuring the natural logarithm is always defined. Furthermore, it acts as mathematical smoothing (similar to Laplace smoothing). It preserves the exact rank-order of the competitors while compressing the artificially inflated log-ratios back to a scale representative of true human skill differences.
	
	We only scale up from below to solve mathematical issues: $T_{floor}$ corrects undefined-log problems and smooths extreme volatility swings on short stages. Scaling long stages down to a reference time introduces issues rather than correcting them:
	
	\begin{enumerate}
		\item \textbf{The ``Speed Tax'' on Field Courses (Reversed Volatility):} 
		While global scaling fixes the volatility of fixed-time mistakes, it introduces massive mathematical volatility for proportional speed differences. If a competitor's physiological pace is exactly 10\% slower than the stage winner, their absolute time difference ($\Delta T$) scales with the length of the stage.
		\begin{itemize}
			\item On a short 10-second stage: $\Delta T = 1.0$ second. Under a universal reference (e.g., $T_{ref} = 20.0$), their ratio is $\frac{21.0}{20.0} = 1.05$ (a 5\% penalty).
			\item On a long 50-second stage: $\Delta T = 5.0$ seconds. Under the same universal reference, their ratio is $\frac{25.0}{20.0} = 1.25$ (a 25\% penalty).
		\end{itemize}
		By standardizing the winner's time universally, the model artificially magnifies the deficit of consistently slower competitors on long field courses.
		
		\item \textbf{Loss of Physical Fidelity:} 
		In practical shooting, the duration of the stage provides critical context for the severity of a mistake. A 3.0-second error on a 5.0-second stage represents a catastrophic failure relative to the task, whereas a 3.0-second error on a 50.0-second stage is a minor inefficiency. A universal additive scale discards the physical reality of the task's duration, treating both errors as mathematically identical.
	\end{enumerate}
	
	Because of these drawbacks, a \textbf{Hybrid Scaling} approach is generally preferred over a universal reference time. By applying the additive padding $k$ only to stages where $T_{min} < T_{floor}$, the model safely bounds the extreme punishment of micro-stages and time bonuses, while leaving standard field courses unadjusted to preserve the mathematical reality of proportional pace.
	
	\section{Benefits Over Glicko-2 and Logistic Models}
	\begin{enumerate}
		\item \textbf{Algorithmic Efficiency:} Glicko-2 requires $\mathcal{O}(N^2)$ pairwise comparisons. A 300-person match yields 44,850 comparisons per competitor. The Latent Log-Ratio model computes the match baseline in $\mathcal{O}(N)$ and adds pairwise corrections against $k$ selected opponents per competitor, yielding $\mathcal{O}(N \cdot k)$ total complexity. With $k = 20$--$30$, a 300-person match requires 6,000--9,000 pairwise comparisons total, rather than 44,850 per competitor.
		
		\item \textbf{Native Representation of Performance Distances:} Logistic models represent the outcome of a pairwise comparison as a probability in $(0, 1)$, which is correct for binary win/loss but introduces information loss for continuous margins. A pairwise expected score of $0.95$ could correspond to a 5\% performance margin or a 50\% margin---the sigmoid saturates at the extremes, compressing large differences into a narrow band. The latent log-ratio model preserves the log-ratio structure of performance distances globally: a competitor who performs at twice the level of another has a log-gap of $\ln(2) \approx 0.69$, regardless of whether both are elite or both are novices. Pathological deep-tail updates are handled by observation-noise inflation rather than by deforming this distance scale, so prediction remains analytically simple.
		
		\item \textbf{Explicit Separation of Variance Sources:} Rather than routing all uncertainty through a single mechanism (e.g., Glicko-2's Rating Deviation, which conflates measurement uncertainty, skill drift, and competitor consistency), the model explicitly separates:
		\begin{itemize}
			\item $\sigma_{sport}^2$: the irreducible noise of the sport for \emph{updates} (a global constant, tuned once),
			\item $\sigma_{pred}^2$: idiosyncratic sport noise used only in \emph{predictions} (win probabilities and predictive bands; typically smaller than $\sigma_{sport}^2$),
			\item $V_i$: measurement uncertainty in each competitor's rating (adapts via the Kalman filter),
			\item $\sigma_i^2$: per-competitor behavioral volatility for \emph{updates} (adapts via EMA); optionally scaled by $\kappa$ for \emph{predictive} bands only,
			\item $\sigma_{drift}^2$: skill drift over time (a global constant for the aging function).
		\end{itemize}
		This decomposition makes each parameter independently interpretable and tunable. In Glicko-2, a competitor bombing a single stage can drastically inflate their RD (via the volatility $\rightarrow$ RD pathway), destabilizing future updates. Here, a single bad performance primarily affects $\sigma_i^2$ (which adapts gradually via the EMA) and may trigger a modest $V_i$ increase (via the innovation term), but neither mechanism produces the runaway instability that Glicko-2's volatility iteration can exhibit.
	\end{enumerate}
	
	\section{The Prediction System and Confidence Intervals}
	\label{sec:prediction}
	
	Because ratings operate in log space, predicted finish ratios use log-normal bookkeeping. The current implementation builds a \textbf{closed-form} mixture over plausible winners; Monte Carlo winner sampling is not wired in (reserved for future contested-field use).
	
	\subsection{Candidate winners and win probabilities}
	Sort competitors by latent rating $R$. Candidate winners are the top $m = \max(5, \lceil 0.25 N\rceil)$ by rating (upper quartile, at least five). For each candidate $w$, a product of pairwise normal CDFs gives an unnormalized weight:
	\begin{equation}
		\tilde{p}_w = \prod_{j \neq w} \Phi\!\left(\frac{R_w - R_j}{\sqrt{V_w + V_j + 2\sigma_{pred}^2}}\right),
	\end{equation}
	where $\Phi$ is the standard normal CDF and $\sigma_{pred}^2$ is the prediction-only sport variance (Section~2). Behavioral volatilities $\sigma_i^2$ are \emph{omitted}: they are not modeled as uncertainty about who wins. Normalize $p_w = \tilde{p}_w / \sum_{w'} \tilde{p}_{w'}$.
	
	\subsection{Expected finish ratio}
	Let $p_w$ be the mixture weights from the previous subsection. By the law of total expectation over the random winner, the expected ratio of competitor $i$ to the winner is
	\begin{equation}
		\hat{S}_i = \sum_w p_w \exp\!\left(R_i - R_w + \frac{V_w + \sigma_{sport}^2}{2}\right).
	\end{equation}
	The term $(V_w + \sigma_{sport}^2)/2$ is the standard log-normal mean correction when the reference (winner) performance is modeled as Gaussian with variance $V_w + \sigma_{sport}^2$ in log space; behavioral volatility does not enter this correction. Posterior variances $V$ are the committed values stored on each rating (not time-aged variance).
	
	\subsection{Predictive variance and reported ``one sigma''}
	Let $\mu^{\mathrm{log}}_i = \sum_w p_w (R_i - R_w)$ be the probability-weighted mean log gap to mixture winners. Let $\bar{\sigma}_w^2 = \sum_w p_w \sigma_w^2$ denote the probability-weighted behavioral variance of the presumed winner mixture. Predictive variance in log space is
	\begin{equation}
		\varsigma_i^2 = V_i + 2\sigma_{pred}^2 + \sum_w p_w V_w + \kappa\left(\sigma_i^2 + \bar{\sigma}_w^2\right).
	\end{equation}
	The sport contribution $2\sigma_{pred}^2$ is idiosyncratic noise for both the focal competitor and the uncertain winner in the ratio. The term $\kappa(\sigma_i^2 + \bar{\sigma}_w^2)$ adds a \emph{fraction} of behavioral variance to the band: $\kappa=0$ recovers bands that ignore $\sigma^2$; the update rule continues to use full $\sigma_i^2$ in $R_{\mathrm{obs},i}$ regardless of $\kappa$.
	
	The UI reports an asymmetric one-standard-deviation \emph{upper} excursion in ratio space: the level $\exp(\mu^{\mathrm{log}}_i + \varsigma_i)$ relative to the mean $\hat{S}_i$,
	\begin{equation}
		\Delta^{\mathrm{1\sigma}}_i = \exp\!\left(\mu^{\mathrm{log}}_i + \varsigma_i\right) - \hat{S}_i,
	\end{equation}
	rather than a symmetric $\pm 1\sigma$ interval or normal 95\% bounds.
	
\end{document}
